\documentclass[11pt]{article}

\usepackage[margin=1in]{geometry}
\usepackage{amsmath,amssymb}
\usepackage{enumitem}
\usepackage{hyperref}
\usepackage{xcolor}
\usepackage{booktabs}
\usepackage{listings}

\hypersetup{colorlinks=true, linkcolor=blue, urlcolor=blue, citecolor=blue}

\lstset{
  basicstyle=\ttfamily\small,
  breaklines=true,
  columns=fullflexible,
}

\title{TOAE209/TOAH209 -- Design and Analysis of Algorithms\\[4pt]
\large Week 1: Practical Exercises}
\author{Foreign Trade University}
\date{}

\begin{document}
\maketitle

\section*{Instructions}

For each problem below there is a starter file
\texttt{exercises/starter\_code\_problemNN.py} containing function signatures with
\texttt{raise NotImplementedError} bodies, and a small \texttt{\_\_main\_\_} block with
tests. Implement the missing functions so the tests pass. A reference solution is
provided in \texttt{solutions/solution\_problemNN.py} -- try the problem yourself before
looking!

All problems use only the Python 3.10+ standard library. Shared helper functions for the
stable-matching problems (Problems 1--8, 18--22) live in \texttt{starter\_code.py} at the
week root: \texttt{generate\_instance}, \texttt{rank}, \texttt{is\_perfect\_matching},
\texttt{find\_rogue\_couples}, \texttt{is\_stable}, \texttt{print\_matching}.

The 22 problems are grouped into three themes:
\begin{itemize}
    \item \textbf{Stable matching} (Problems 1--8, 18--20, 22): the Gale--Shapley
    algorithm and its theoretical properties.
    \item \textbf{Data structures} (Problems 9--12, 16): arrays, linked lists, stacks,
    queues, and heaps.
    \item \textbf{Algorithmic building blocks} (Problems 13--17, 21): recursion, searching,
    sorting, divide-and-conquer multiplication, and asymptotic growth.
\end{itemize}

\newpage
\section{Stable Matching: Core Algorithm}

\subsection*{Problem 1 -- Perfect Matching Checker}

Implement \texttt{is\_perfect\_matching(matching, men, women)} \textbf{from scratch}
(without using the helper of the same name in \texttt{starter\_code.py}). A matching
(a dict \texttt{man -> woman}) is \emph{perfect} iff:
\begin{enumerate}
    \item Every man in \texttt{men} appears exactly once as a key.
    \item Every woman in \texttt{women} appears exactly once as a value, i.e. the matching
    is a bijection between \texttt{men} and \texttt{women}.
\end{enumerate}

\textbf{Example.} With \texttt{men = ["M0","M1","M2"]}, \texttt{women =
["W0","W1","W2"]}, the matching \texttt{\{"M0":"W1","M1":"W0","M2":"W2"\}} is perfect, but
\texttt{\{"M0":"W1","M1":"W1","M2":"W2"\}} is not (W1 used twice, W0 missing).

\subsection*{Problem 2 -- Stability Checker}

Implement \texttt{find\_rogue\_couples(matching, men\_prefs, women\_prefs)}
\textbf{from scratch}. A pair $(m, w)$ is a \emph{rogue couple} (blocking pair) for
\texttt{matching} if:
\begin{itemize}
    \item $m$ and $w$ are \emph{not} matched to each other, \textbf{and}
    \item $m$ prefers $w$ to his current partner, \textbf{and}
    \item $w$ prefers $m$ to her current partner.
\end{itemize}
Return a list of all such $(m,w)$ pairs. A matching is \emph{stable} iff this list is
empty.

\textbf{Example.} Two men $\{M0, M1\}$, two women $\{W0, W1\}$, where both men prefer W0
and both women prefer M0 (i.e. everyone agrees W0 > W1 and M0 > M1). The matching
$\{M0{:}W0, M1{:}W1\}$ is stable. The matching $\{M0{:}W1, M1{:}W0\}$ is \emph{unstable}:
$(M0, W0)$ is a rogue couple, since $M0$ prefers $W0$ to $W1$ and $W0$ prefers $M0$ to
$M1$.

\subsection*{Problem 3 -- Gale--Shapley Algorithm (Men Propose)}

Implement \texttt{gale\_shapley\_men\_propose(men\_prefs, women\_prefs)}:
\begin{verbatim}
Initialize all m in M and w in W to free.
While some man m is free and hasn't proposed to every woman:
    w = first woman on m's list to whom m has not yet proposed
    if w is free:
        (m, w) become engaged
    else if w prefers m to her current partner m':
        m' becomes free; (m, w) become engaged
    else:
        w rejects m
\end{verbatim}
Return the resulting matching as a dict \texttt{man -> woman}. Verify on a random
instance (via \texttt{generate\_instance}) that the result is a perfect matching
(Problem 1) and stable (Problem 2).

\subsection*{Problem 4 -- Gale--Shapley Algorithm (Women Propose)}

By symmetry with Problem 3, implement \texttt{gale\_shapley\_women\_propose(men\_prefs,
women\_prefs)}, where the roles are swapped: women propose to men, and men hold the best
offer received so far. Return the matching in the \emph{same orientation} as Problem 3
(\texttt{man -> woman}), so the two results are directly comparable. \emph{Hint}: you can
implement this by calling \texttt{gale\_shapley\_men\_propose} with the roles of
\texttt{men\_prefs} and \texttt{women\_prefs} swapped, then flipping the resulting dict.

\subsection*{Problem 5 -- Brute-Force Enumeration of All Stable Matchings}

For small $n$, implement \texttt{all\_stable\_matchings(men, women, men\_prefs,
women\_prefs)} that returns a list of \emph{all} stable perfect matchings, by:
\begin{enumerate}
    \item Generating every possible perfect matching (every permutation of \texttt{women}
    assigned to \texttt{men}, i.e. $n!$ candidates), using
    \texttt{itertools.permutations}.
    \item Keeping only those that are stable (Problem 2 / \texttt{is\_stable}).
\end{enumerate}
This is exponential and only intended for small $n$ (say $n \le 6$); it exists so we can
empirically verify properties of Gale--Shapley in Problems 6--8.

\subsection*{Problem 6 -- How Many Stable Matchings Are There?}

Using \texttt{all\_stable\_matchings} from Problem 5, implement
\texttt{average\_num\_stable\_matchings(n, num\_trials, seed)}:
\begin{enumerate}
    \item Generate \texttt{num\_trials} random instances of size $n$ (with
    \texttt{generate\_instance(n, seed=seed + trial\_index)} or similar, so each trial is
    different but reproducible).
    \item Compute the number of stable matchings for each instance.
    \item Return the average over all trials, as a float.
\end{enumerate}
In \texttt{\_\_main\_\_}, print this average for $n = 1, \dots, 5$ with
\texttt{num\_trials=20} and observe how quickly the count grows with $n$.

\subsection*{Problem 7 -- Verifying the Men-Optimality Theorem}

\textbf{Theorem (Kleinberg \& Tardos, Thm 1.7).} In the matching produced by the
men-proposing Gale--Shapley algorithm, every man receives the \emph{best valid partner}
he could have in \emph{any} stable matching (he is ``man-optimal'').

A woman $w$ is a \emph{valid partner} of man $m$ if there exists \emph{some} stable
matching in which $m$ and $w$ are matched.

Implement:
\begin{itemize}
    \item \texttt{valid\_partners(person, all\_stable, role)} -- returns the set of valid
    partners of \texttt{person}, where \texttt{role} is \texttt{"man"} or
    \texttt{"woman"}, using the list of all stable matchings from Problem 5.
    \item \texttt{verify\_men\_optimality(men, men\_prefs, all\_stable, men\_optimal)} --
    checks, for every man $m$, that \texttt{men\_optimal[m]} is $m$'s
    \emph{most-preferred} valid partner.
\end{itemize}
Test by generating a random instance, computing \texttt{all\_stable\_matchings} (Problem
5) and \texttt{gale\_shapley\_men\_propose} (Problem 3), and confirming
\texttt{verify\_men\_optimality} returns \texttt{True}.

\subsection*{Problem 8 -- Verifying the Women-Pessimality Theorem}

\textbf{Theorem (Kleinberg \& Tardos, Thm 1.8).} In the same matching, every woman is
matched with her \emph{least-preferred} valid partner (she is ``woman-pessimal'').

Using \texttt{valid\_partners} from Problem 7, implement
\texttt{verify\_women\_pessimality(women, women\_prefs, all\_stable, men\_optimal)} that
checks, for every woman $w$, that her partner in \texttt{men\_optimal} is her
\emph{least}-preferred valid partner. Test the same way as Problem 7.

\newpage
\section{Data Structures}

\subsection*{Problem 9 -- Array vs. Linked List: Counting Operations}

Implement two simple sequence data structures, each instrumented with an \texttt{ops}
counter that counts ``basic steps'' (element shifts for the array, pointer hops for the
linked list):

\begin{itemize}
    \item \texttt{ArraySeq.prepend(x)}: insert \texttt{x} at index 0. Every existing
    element must be shifted right by one slot $\to$ costs \texttt{len(self.data)} ops.
    \item \texttt{ArraySeq.get(i)}: direct indexing $\to$ costs 1 op.
    \item \texttt{LinkedSeq.prepend(x)}: create a new head node $\to$ costs 1 op.
    \item \texttt{LinkedSeq.get(i)}: walk \texttt{i} pointers from the head $\to$ costs
    $(i+1)$ ops.
\end{itemize}

Then implement \texttt{total\_prepend\_ops(seq\_class, n)} that creates a fresh
\texttt{seq\_class()}, prepends $n$ items ($0, 1, \dots, n-1$ in that order), and returns
the final value of \texttt{seq.ops} (the total cost of all prepends).

\textbf{What you should observe.} For \texttt{ArraySeq}, the total is
$0+1+\cdots+(n-1) = \binom{n}{2} = \Theta(n^2)$. For \texttt{LinkedSeq}, the total is
exactly $n = \Theta(n)$. This is a direct, measurable illustration of the array-vs-linked-list
trade-off from the lecture notes.

\subsection*{Problem 10 -- Stack: Balanced Brackets Checker}

Implement a simple \texttt{Stack} class with \texttt{push}, \texttt{pop}, \texttt{peek},
and \texttt{is\_empty}, backed by a Python list. Then implement \texttt{is\_balanced(expr)}
that returns \texttt{True} iff every \texttt{(}, \texttt{[}, \texttt{\{} in \texttt{expr}
is closed by the matching \texttt{)}, \texttt{]}, \texttt{\}} in the correct order (other
characters are ignored).

\textbf{Examples.} \texttt{"(a[b]\{c\})"} $\to$ \texttt{True}. \texttt{"(a[b)]"} $\to$
\texttt{False} (wrong order). \texttt{"((a)"} $\to$ \texttt{False} (unclosed).

\subsection*{Problem 11 -- Queue: Round-Robin CPU Scheduling}

Implement a \texttt{Queue} class backed by a Python list with \texttt{enqueue},
\texttt{dequeue}, and \texttt{is\_empty} (FIFO order).

Then implement \texttt{round\_robin(tasks, quantum)}, where \texttt{tasks} is a list of
\texttt{(name, burst\_time)} pairs. Simulate round-robin scheduling: repeatedly dequeue a
task, run it for up to \texttt{quantum} time units, and if it still has remaining time,
enqueue it again at the back. Return the order in which tasks \emph{finish} (a list of
names, in finishing order).

\textbf{Example.} \texttt{tasks = [("A",5), ("B",3), ("C",8)]}, \texttt{quantum=4}.
Trace: A runs 4 (1 left), B runs 3 (done), C runs 4 (4 left), A runs 1 (done), C runs 4
(done). Finishing order: \texttt{["B", "A", "C"]}.

\subsection*{Problem 12 -- Singly Linked List}

Implement a \texttt{LinkedList} class with:
\begin{itemize}
    \item \texttt{append(x)}: add \texttt{x} at the end -- $O(n)$ without a tail pointer.
    \item \texttt{prepend(x)}: add \texttt{x} at the front -- $O(1)$.
    \item \texttt{delete(x)}: remove the first node with value \texttt{x} (no-op if not
    found).
    \item \texttt{find(x)}: return \texttt{True} iff \texttt{x} is in the list.
    \item \texttt{to\_list()}: return a Python list of the values, in order.
    \item \texttt{reverse()}: reverse the list \emph{in place} (re-link nodes -- do not
    just reverse \texttt{to\_list()}).
\end{itemize}

\subsection*{Problem 16 -- Binary Min-Heap and Heapsort}

Implement an array-based binary \texttt{MinHeap} with:
\begin{itemize}
    \item \texttt{push(x)}: insert \texttt{x}, then ``sift up'' to restore the heap
    property.
    \item \texttt{pop()}: remove and return the minimum element, moving the last element
    to the root and ``sifting down'' to restore the heap property. Raise
    \texttt{IndexError} if empty.
    \item \texttt{\_\_len\_\_}: number of elements.
\end{itemize}
Then implement \texttt{heapsort(arr)} that returns a new sorted list by pushing all
elements onto a \texttt{MinHeap} and popping them off in order. This gives an
$O(n\log n)$ sort.

\newpage
\section{Algorithmic Building Blocks}

\subsection*{Problem 13 -- Binary Search (Iterative and Recursive)}

Implement two versions of binary search over a sorted list \texttt{arr}:
\begin{itemize}
    \item \texttt{binary\_search\_iterative(arr, target)}: returns the index of
    \texttt{target} in \texttt{arr}, or $-1$ if not present. Use a \texttt{while} loop.
    \item \texttt{binary\_search\_recursive(arr, target)}: same behaviour, implemented
    recursively (a helper with \texttt{lo}/\texttt{hi} bounds is fine).
\end{itemize}
Also implement \texttt{recursion\_depth(n)}, returning the number of recursive calls
\texttt{binary\_search\_recursive} would make in the \emph{worst case} on a sorted array
of length $n$, i.e. $\lceil \log_2(n+1) \rceil$ -- the number of halvings until the search
range is empty.

\subsection*{Problem 14 -- Insertion Sort with Comparison Counter}

Implement \texttt{insertion\_sort(arr)} that returns a \emph{new} sorted list (do not
mutate the input) using the standard insertion-sort algorithm, together with the total
number of element \emph{comparisons} performed, as a tuple
\texttt{(sorted\_list, comparisons)}.

\textbf{Verify empirically} that:
\begin{itemize}
    \item On an already-sorted list of length $n$, insertion sort makes exactly $n-1$
    comparisons (best case, $O(n)$).
    \item On a reverse-sorted list of length $n$, it makes exactly $\frac{n(n-1)}{2}$
    comparisons (worst case, $O(n^2)$).
\end{itemize}

\subsection*{Problem 15 -- Karatsuba Multiplication}

Implement \texttt{karatsuba(x, y)}, multiplying two non-negative integers using the
Karatsuba divide-and-conquer algorithm: given $x, y$ with $m$ digits each, split
$x = a\cdot 10^{m/2} + b$ and $y = c\cdot 10^{m/2} + d$. Then
\[
  ac = a \cdot c, \qquad bd = b \cdot d, \qquad
  ad\_plus\_bc = (a+b)(c+d) - ac - bd,
\]
\[
  x \cdot y = ac \cdot 10^{m} + ad\_plus\_bc \cdot 10^{m/2} + bd
\]
using \emph{three} recursive multiplications instead of four. Base case: if $x < 10$ or
$y < 10$, return $x \cdot y$ directly.

Also implement \texttt{naive\_multiply\_op\_count(x, y)} that returns the number of
single-digit multiplications a grade-school algorithm would need, roughly
\texttt{len(str(x)) * len(str(y))}.

\subsection*{Problem 17 -- Asymptotic Growth Classification}

You are given a dictionary mapping function names to callables of $n$:
\begin{verbatim}
FUNCTIONS = {
    "constant":     lambda n: 1,
    "logarithmic":  lambda n: math.log2(n),
    "sqrt":         lambda n: math.sqrt(n),
    "linear":       lambda n: n,
    "linearithmic": lambda n: n * math.log2(n),
    "quadratic":    lambda n: n ** 2,
    "cubic":        lambda n: n ** 3,
    "exponential":  lambda n: 2 ** n,
}
\end{verbatim}
Implement \texttt{order\_by\_growth(functions, n)} that returns the list of function names
sorted by \emph{increasing} value at the given $n$ (i.e. from slowest- to
fastest-growing at that $n$).

Then implement \texttt{dominates(f, g, n\_values)} that returns \texttt{True} iff, for
every $n$ in \texttt{n\_values} (a list of increasing positive integers),
$f(n) \le g(n)$ -- a simple empirical proxy for ``$f$ is $O(g)$''.

\subsection*{Problem 21 -- Recursion: Naive vs. Memoized Fibonacci}

Implement \texttt{fib\_naive(n)} returning \texttt{(value, call\_count)} where
\texttt{value} is the $n$-th Fibonacci number ($F(0)=0$, $F(1)=1$) and
\texttt{call\_count} is the total number of recursive calls made (including the initial
call).

Then implement \texttt{fib\_memo(n)} returning \texttt{(value, call\_count)} using
memoization (a dict cache), where \texttt{call\_count} should be much smaller -- $O(n)$
instead of exponential.

\textbf{What you should observe.} For $n=20$: \texttt{fib\_naive} makes 21891 calls, while
\texttt{fib\_memo} makes only 39.

\newpage
\section{Stable Matching: Extensions \& Trace Practice}

\subsection*{Problem 18 -- Empirical Running Time of Gale--Shapley}

Implement \texttt{count\_proposals(men\_prefs, women\_prefs)}, a variant of the
men-proposing Gale--Shapley algorithm (Problem 3) that returns the \emph{total} number of
proposals made (instead of the final matching).

The Kleinberg--Tardos textbook proves this is at most $n^2$ (Section 1.1's analysis).
Implement \texttt{average\_proposals(n, num\_trials, seed)} returning the average number
of proposals over \texttt{num\_trials} random instances of size $n$, and verify it never
exceeds $n^2$.

\subsection*{Problem 19 -- Hospital--Residents Problem (Many-to-One Stable Matching)}

Generalize Gale--Shapley to the Hospital--Residents problem: each hospital $h$ has a
capacity \texttt{capacity[h]} (it can accept up to that many residents). Residents propose
to hospitals; a hospital holds the best \texttt{capacity[h]} residents it has seen so far
(by its preference list) and rejects the rest.

Implement \texttt{hospital\_resident\_matching(residents, hospitals, resident\_prefs,
hospital\_prefs, capacity)}:
\begin{itemize}
    \item \texttt{residents}, \texttt{hospitals}: lists of names.
    \item \texttt{resident\_prefs[r]}: hospitals in order of preference (best first).
    \item \texttt{hospital\_prefs[h]}: residents in order of preference (best first).
    \item \texttt{capacity[h]}: max number of residents hospital $h$ can take.
\end{itemize}
Return a dict \texttt{hospital -> list of assigned residents}, where every hospital's list
has length $\le$ \texttt{capacity[h]}, and the assignment is stable: no resident $r$ and
hospital $h$ both prefer each other over $r$'s current assignment / one of $h$'s current
residents.

\subsection*{Problem 20 -- Stable Matching with Incomplete Preference Lists}

In practice, not everyone is acceptable to everyone else. Generalize Gale--Shapley to
\emph{incomplete} preference lists: \texttt{men\_prefs[m]} and \texttt{women\_prefs[w]}
list only the \emph{acceptable} partners, best first (a person may have an empty list).

Run the men-proposing algorithm: each man proposes down his list. If he runs out of
acceptable women, he remains permanently single. Women only accept proposals from men on
their own list.

Implement \texttt{gale\_shapley\_incomplete(men, women, men\_prefs, women\_prefs)}
returning a dict \texttt{man -> woman} containing \emph{only} the matched pairs
(unmatched men/women are simply absent from the dict).

\subsection*{Problem 22 -- Trace-Based Verification on a Textbook Example}

Consider the classic $3\times3$ instance (Kleinberg \& Tardos, Chapter 1):

\begin{center}
\begin{tabular}{@{}ll@{}}
\toprule
Men's preferences (best first) & Women's preferences (best first) \\
\midrule
$X$: A, B, C & $A$: Y, X, Z \\
$Y$: B, A, C & $B$: X, Y, Z \\
$Z$: A, B, C & $C$: X, Y, Z \\
\bottomrule
\end{tabular}
\end{center}

Implement \texttt{build\_instance()} returning \texttt{(men, women, men\_prefs,
women\_prefs)} for this instance, where \texttt{men = ["X","Y","Z"]} and
\texttt{women = ["A","B","C"]}.

Using \texttt{gale\_shapley\_men\_propose} from Problem 3, the men-proposing algorithm on
this instance produces the matching $\{X{:}A, Y{:}B, Z{:}C\}$. You can verify this by
hand-tracing the algorithm:

\begin{center}
\begin{tabular}{@{}clll@{}}
\toprule
Step & Proposal & Outcome & Engaged pairs after step \\
\midrule
1 & $X \to A$ & $A$ free, accepts & $(X,A)$ \\
2 & $Y \to B$ & $B$ free, accepts & $(X,A), (Y,B)$ \\
3 & $Z \to A$ & $A$ prefers $X$ to $Z$ (A's list: Y,X,Z) -- $Z$ rejected & $(X,A), (Y,B)$ \\
4 & $Z \to B$ & $B$ prefers $Y$ to $Z$ (B's list: X,Y,Z) -- $Z$ rejected & $(X,A), (Y,B)$ \\
5 & $Z \to C$ & $C$ free, accepts & $(X,A), (Y,B), (Z,C)$ \\
\bottomrule
\end{tabular}
\end{center}

Note that $Z$, the only man rejected along the way, ends up with his \emph{last} choice
$C$ -- while $A$ and $B$ (the women he was rejected by) each end up with their
\emph{first}-choice man among those who proposed to them. This single trace illustrates
both the men-optimal and women-pessimal theorems from Problems 7--8.

Implement \texttt{verify\_textbook\_example()} that builds the instance, runs
Gale--Shapley, and returns \texttt{True} iff the result equals
\texttt{\{"X":"A","Y":"B","Z":"C"\}}.

\end{document}
